\section{Introduction}
\label{sec:intro}

A video world model for manipulation takes the current camera frame and a sequence of robot actions and predicts the video the robot would record.
If the prediction is right, it can replace a physics simulator for policy evaluation, planning and data generation \cite{zhu2024irasim,nvidia2025cosmos}.
These models are usually compared with frame metrics (PSNR, SSIM \cite{wang2004ssim}, LPIPS \cite{zhang2018lpips}) or distribution metrics such as FVD.
Such scores reward a video that looks like the recorded one; they do not ask whether the robot in it executes the given actions, whether objects move only when touched, or whether the task ends as it should.

This report studies that gap on real robot data. Our contributions:
\begin{itemize}
\item a reproduction of IRASim and its three released baselines on Bridge \cite{walke2023bridgedata} and Language-Table \cite{lynch2022interactive}, with paired per-clip statistics (Sec.~\ref{sec:reproduce});
\item an evaluation framework with five metric groups, each stated with what it measures and what it cannot establish, applied to six models and two references (Sec.~\ref{sec:framework}, \ref{sec:results});
\item a perturbation test showing that pixel metrics prefer a frozen video to a generation from corrupted actions, while trajectory scores do not (Sec.~\ref{sec:results});
\item a reference-free 3-D action-consistency score built on an off-the-shelf inverse dynamics model, validated on real clips, and its use for best-of-N selection and for an executability check in simulation (Sec.~\ref{sec:improve}).
\end{itemize}
